\section{P: independent splices and the changed-vertex barrier}\label{sec:splice}
S Section 3 uses only square-lattice rotations and reflections for its charts.
Fix one chart with upper labels having contact directions strictly above
the horizontal axis and lower labels strictly below it. These are D4
isometries, so coordinate Lipschitz constants and nearest-neighbor height
counts are unchanged. The diagonal wall coordinates used earlier are a
separate construction. The angular gap may depend on a label; no uniform
gap over the label set is asserted.

S Section 4 transfers exceptional labels by independent half-plane pasting.
Its exact shared-weight premise fails for caps. We now prove its required
transfer conclusion with actual changed vertices in place of the cut.
Let $E_n^+$ contain edges whose two endpoint heights are at least $-n$.
Form $X_n^+$ from $X$ on $E_n^+$ and an independent field elsewhere; form
$X_n^-$ from another independent field on $E_n^+$ and $X$ elsewhere.
Splice both originals and cap auxiliaries. For each fixed $n$, the two outputs
are independent and each has the original product law. Their caps $s_n^+,s_n^-$
are therefore independent with the cap law. Independence across $n$ is not needed.

Let $D_n^+$ be all endpoints of edges where $s$ and $s_n^+$ differ, and
define $D_n^-$ similarly. If an unchanged original edge has protecting
minimum circuits in both fields entirely inside the shared region, the
two minima are mutually admissible and its caps agree. Thus, with $|e|$
the midpoint norm, every changed vertex satisfies
\begin{align}
 z\in D_n^+&\Longrightarrow z_2\le-n+3+
       \max_{e\ni z}\big(L_e(X)+L_e(X_n^+)\big),\label{eq:band}\\
 z\in D_n^-&\Longrightarrow z_2\ge-n-3-
       \max_{e\ni z}\big(L_e(X)+L_e(X_n^-)\big).\nonumber
\end{align}
The constants cover endpoint and midpoint conventions. These sets can be
infinite and need not be curves or connected.

Define
\[
 A_n(X)=\sup_e\frac{L_e(X)}{|e|+n},\qquad
 E_n(T)=\sup_{x,y}\frac{|T(x,y)-\mu(y-x)|}{|x|+|y|+n}.
\]
Uniform sublinear coding radii and \eqref{eq:extended} imply that both tend
to zero almost surely. The same quantities for $X_n^+$ tend to zero in
probability, since each has the same marginal law. Choose one deterministic
subsequence $n_k$ on which their probabilities of exceeding $2^{-k}$ are
summable. Borel--Cantelli gives, on that subsequence, all four errors tending
to zero. Consequently for random $\eps_k\downarrow0$,
\begin{equation}\label{eq:movingband}
 z\in D_{n_k}^+\quad\Longrightarrow\quad
 z_2\le-n_k+\eps_k(|z|+n_k).
\end{equation}
For every fixed $n$, both fields' coding radii also give
$z_2\le o(|z|)+O_n(1)$ on $D_n^+$ and the reflected bound on $D_n^-$,
simultaneously for the countable collection of $n$.

\begin{lemma}[Moving and fixed barriers]\label{lem:splice}
Every finite collection of upper-ray Busemann witnesses in $s$ transfers
to $s_{n_k}^+$ for all sufficiently large $k$, with the same exact labels.
For every fixed $n$, every lower ray has a tail in $s_n^-$ with its exact
original label. These statements hold pathwise for all rays on the
structural full event.
\end{lemma}
\begin{proof}
Fix an upper ray $p$, its functional $\rho$, and a finite set of basepoints
$K$ containing $p_0$. Compactness and the strict upper contact sector give
$\eta,\gamma>0$ such that
\begin{equation}\label{eq:gap}
 \mu(d)-\rho(d)\ge\gamma|d|\quad\text{if }d_2\le\eta|d|.
\end{equation}
For large $k$, \eqref{eq:movingband} implies $|z|\ge n_k/2$ for all
$z\in D_{n_k}^+$. Uniformly for $x\in K$, $|z-x|\ge|z|/2$ and
$|x|+|z|+n_k\le8|z-x|$. Also
$(z-x)_2\le3\eps_k|z|+O_K(1)$, so \eqref{eq:gap} applies.
The two metric errors and the Busemann error are small compared with
$|z-x|$. Hence, uniformly on the infinite changed set,
\begin{equation}\label{eq:movingbarrier}
 \min(T,T_n^+)(x,z)-B_p(x,z)
 \ge(\gamma/2)|z-x|\ge\gamma n_k/8.
\end{equation}
The ray itself avoids $D_{n_k}^+$ for large $k$: its eventual strict upper
sector excludes distant vertices by \eqref{eq:movingband}, and the receding
cut excludes its finite prefix.

For a walk from $x$ to $p_j$ meeting $D_{n_k}^+$, let $z$ be its \emph{last}
vertex in that set. Its suffix has exactly the same edge costs in both
fields. Calibration and the metric bound for $B_p$ show that the new walk
costs at least $T_n^+(x,z)+B_p(z,p_j)$, which exceeds
$B_p(x,p_j)$ by the barrier. The same bound holds in the old field. Walks
avoiding the changed set have exactly equal costs. With $x=p_0$ this preserves
all prefixes of $p$. For any fixed $x\in K$ the calibration defect
$T(x,p_j)-B_p(x,p_j)$ tends to zero, so eventually both minimizers avoid
the changed set and
\begin{equation}\label{eq:distanceshare}
 T_n^+(x,p_j)=T(x,p_j).
\end{equation}
This preserves the requested finite Busemann witnesses. The subsequence
was chosen before $p$, so the argument is simultaneous over rays without
new probability intersections.

To identify the exact label fix $n$ and a preserved upper ray. For any
$c>0$, take $x_2\ge c|x|$. The fixed-cut bound gives
$z_2\le\delta|z|+C_\delta$ on $D_n^+$. For small $\delta$ and large
$|x|$, elementary triangle inequalities yield
\[
 |z-x|\ge\kappa(|x|+|z|),\qquad \kappa=c/(c+8).
\]
Indeed $|z-x|<c|x|/4$ would imply $|z|<2|x|$ and
$x_2-z_2>c|x|/2$, a contradiction when $\delta\le c/8$.
Decrease $\delta$ so $\delta/\kappa<\eta/2$. The displacement lies in
the cone of \eqref{eq:gap}; the fixed-metric and Busemann errors then give
\[
 \min(T,T_n^+)(x,z)-B_p(x,z)
 \ge(\gamma\kappa/4)(|x|+|z|)
\]
for all sufficiently far upper-sector $x$, uniformly over $z\in D_n^+$.
The last-changed-vertex argument gives equality of Busemann values at those
$x$. Use $x=ke_2$ and $x=k(e_1+e_2)$ to identify the linear functional
in the new field as exactly $\rho$. The reflected argument applies for
each fixed $n$ to a sufficiently late tail of every lower ray in $s_n^-$,
using two independent lower directions. Thresholds may depend on the ray
and on $n$, as in S's transfer statement.
\end{proof}

S Section 4's remaining proof uses the countability of exceptional labels,
Borel/analytic witness sets and independence of the two output fields.
Lemma~\ref{lem:splice} supplies exactly its two transfer quantifiers,
including finite witnesses and exact labels; fixed-label coalescence is
available from Section~\ref{sec:allray}. Its four-witness independence
argument therefore gives opposite exceptional coalescence simultaneously.
No exact half-plane agreement of capped weights and no original Busemann
linearity have been assumed.
